% Corrigé intégré automatiquement pour la banque Exercices.
% Source : A_Integrer/DDS_04/078_Modelisation/078_Modelisation.tex
\begin{corrigebox}[Corrigé]
\CorrigeQuestion{1}
$F=2F_1=\SI{357}{N}$.
\CorrigeQuestion{2}
$F=\dfrac{\mu_0 N^2 I^2 S}{8\left(e_0 - x  \right)^2}$, $\gamma = \SI{1,9e-6}{Nm^2A^{-2}}$.
\CorrigeQuestion{3}
La résolution consiste à identifier les données et l'inconnue, choisir la loi du cours adaptée, écrire l'équation symbolique avant toute application numérique, puis vérifier l'homogénéité, le signe et l'ordre de grandeur du résultat. Les hypothèses utilisées doivent être explicitement contrôlées à la fin.
\CorrigeQuestion{4}
On choisit le système isolé et le référentiel galiléen, puis on calcule le torseur cinétique et le torseur dynamique au point le plus commode. Le PFD s'écrit $\{\mathcal T_{ext}\}=\{\mathcal D\}$ ; ses équations scalaires donnent les efforts ou la loi de mouvement.
\CorrigeQuestion{5}
$H_{\text{PM}}=\dfrac{8e_0\gamma k_a^2 U_0}{mp^2e_0^3-8\gamma k_a^2 U_0^2}$.
\end{corrigebox}
